% !PRINTED 127
\weberpage{127}
%<<CONT
wozu noch folgendes zu bemerken ist: Die Wurzeln $\sqrt{\alpha+\beta\omega}$,
$\sqrt{-i(\alpha+\beta\omega)}$ sind mit \gesperrt{positivem} reellem Teil zu nehmen.
Daß eine von ihnen rein imaginär sei, ist durch die Annahme, daß $\omega$
einen positiven imaginären Teil hat und $\alpha$ bzw. $\beta$ positiv sei,
ausgeschlossen, denn danach kann $\alpha+\beta\omega$ oder $-i(\alpha+\beta\omega)$
nicht reell und negativ sein; $\left(\dfrac{\beta}{\alpha}\right)$ und
$\left(\dfrac{\alpha}{\beta}\right)$ ist das \gesperrt{Legendre-Jacobische} Symbol
aus der Theorie der quadratischen Reste, mit der Erweiterung, daß
$\left(\dfrac{\beta}{1}\right)$ und $\left(\dfrac{0}{1}\right)=1$ sein soll. Wenn
im ersten Falle $\alpha$ oder im zweiten $\beta$ negativ ist, so müssen rechts
die sämtlichen Vorzeichen von $\alpha, \beta, \gamma, \delta$ umgekehrt werden.
Wenn sowohl $\alpha$ als $\beta$ ungerade sind, so kann sowohl (15), 1. als (15), 2.
angewandt werden, und beides ergibt, wie man leicht auf Grund
des Reziprozitätsgesetzes der quadratischen Reste nachweist, dasselbe
 Resultat. Es ist nämlich, wenn $\alpha$ und $\beta$ ungerade sind,
$\alpha$ positiv angenommen wird, und das obere oder untere Zeichen
gilt, je nachdem $\beta$ positiv oder negativ ist:
\[
\left(\frac{\beta}{\alpha}\right)\left(\frac{\pm\alpha}{\pm\beta}\right)
= \pm\, e^{-\frac{\pi i}{4}(\alpha\mp1)(\beta-1)},
\]
\[
\sqrt{\mp i(\alpha+\beta\omega)} = e^{\mp\frac{\pi i}{4}}\sqrt{\alpha+\beta\omega}
\]
und die Identität von (15), 1., 2. ergibt sich dann aus den Kongruenzen

\[
\begin{gathered}
-3\alpha\beta+\beta(\alpha+\delta)-(\beta^2-1)\alpha\gamma \equiv \alpha(\gamma-\beta)-(\alpha^2-1)\beta\delta \ (\mathrm{mod}\ 24), \\
\alpha\gamma+\alpha\beta \equiv \beta\delta \ (\mathrm{mod}\ 8),
\end{gathered}
\]
von denen die zweite, wenn $\alpha$ und $\beta$ beide nicht durch 3 teilbar,
also $\alpha^2\equiv\beta^2\equiv1 \ (\mathrm{mod}\ 3)$ sind, auch für den Modul 24 besteht.

Daß durch (15) die Formeln (8) bis (11) befriedigt sind, ist
unmittelbar einzusehen, und es bleibt noch zu zeigen, daß (13)
und (14) erfüllt sind. Wir beginnen mit (14), wobei angenommen
werden kann, daß $\alpha$ ungerade (und positiv) sei; denn vertauscht
man in (14) $\omega$ mit $-1:\omega$, so vertauschen sich $\alpha$ und $-\beta$, und
diese können nicht beide gerade sein.

Es ist
\[
\sqrt{-i\omega}\ \sqrt{\alpha-\frac{\beta}{\omega}} = \sqrt{-i(-\beta+\alpha\omega)};
\]
denn setzen wir für den Augenblick
%CONT>>
